\scanpage{31}
% @ 1
\uline{Prop.} $G$ \deleted{\unsure{groupe alg.}} \inserted{groupe \inserted{(monoïde)}} affine sur le corps $k$. Alors

% @ 2–3
a) $G$ \deleted{\unsure{G}} est limite projective \inserted{\unsure{filtrante|filtr. de|filtrants}} de \inserted{groupes} (monoïdes) affines de t.f. sur $k$, \deleted{\unsure{et}} \deleted{\unsure{b} avec des morphismes}

% @ 4–5
b) Si $G$ est un monoïde \inserted{(groupe)} affine de t.f. sur $k$, on peut trouver un \deleted{\unsure{i|imm.}} entier $n$ et une immersion fermée
% @ 6
\[
 G\unsure{-|\to}\underset{\text{\inserted{schéma \unsure{de|des} matrices $n\times n$}}}{\deleted{\unsure{\mathrm{GL}}}M(n)_{k}^{\unsure{\times|x}}}.\quad
 (\text{resp. }G\to\mathrm{GL}(n))
\]

% @ 7–9
\uline{Dém} $G\longleftrightarrow A$, \deleted{\unsure{et}} et $A=\underset{i}{\unsure{\underrightarrow{\lim}|\underline{\lim}}}\,A_i$, $A_i$ les sous-algèbres de t. f. de $A$. On s'intéresse aux sous-algèbres $B$ qui sont ``stables'' par $\pi:A\otimes A$, or il est clair que si
% @ 10–12
$M\subset A$ est stable sous $\pi$, alors aussi la sous-algèbre \inserted{$B$} qu'elle engendre. Donc le corollaire 2 nous \unsure{dit|est} que l'on peut supposer les $A_i$ stables par $\pi$. Si $G$ un groupe, d'où un \unsure{antipodisme|antipodismes}
% @ 13–17
$\sigma$ de $A$, alors si $M$ stable sous $\sigma$, alors \unsure{l'algèbre} $B$ aussi. Or si $M$ stable par $\pi$, $\sigma M$ aussi [$M$ sous-module à g. \inserted{\unsure{(à dr)|(à d.)}} sur $\check A$ $\Longleftrightarrow$ $\sigma M$ sous-module à dr (g) sur $\check A$] donc $M+\sigma M$ aussi, donc on peut supposer les $A_i$ \unsure{de plus|d'ailleurs} \unsure{stables|stable} par $\sigma$, donc les $G_i$ sont des groupes.

% @ 18–21
Supposons \unsure{maintnt|maintenant} $A$ de t-f., donc on a un sous $k$-module \inserted{de t.f.} $M$ de $A$ qui engendre $A$ comme algèbre. \deleted{Considérons} On peut supposer $M$ stable par $\check A$ à \deleted{\unsure{g.}} \deleted{\unsure{droite}} \inserted{droite} donc que $G$ y opère à \deleted{\unsure{gauche}} \inserted{droite}. Donc
% @ 22
\[
 G\longrightarrow\deleted{\unsure{\mathrm{GL}|{\mathrm{End}}}}\,\underline{\mathrm{Aut}}_{k}(M)\quad(\simeq\mathrm{GL}(n)_{k})
\]
% @ 23–24
\unsure{Av|où|av} $n=\dim M$. Prouvons que c'est un mono. \deleted{\unsure{On a}} Soit \deleted{\unsure{g}\ \unsure{\varepsilon(f)}\ \unsure{=0}} Considérons \inserted{pour \unsure{tt} $\varphi\in M$} la fonction sur $G$
% @ 25
\[
 G\longrightarrow\unsure{\mathbb E|E}\qquad \deleted{\unsure{g}}\quad
 g\longmapsto\varepsilon(g.\varphi)=(g\varphi)(1)=\varphi(g)
\]
% @ 26–28
\deleted{\unsure{donc} $E'$} Donc \deleted{\unsure{g}} $M\subset\mathrm{Im}\,[\text{anneau }\inserted{\text{affine}}\text{ de }\underline{\mathrm{End}}_k(M)\mathrel{-}A]$ \deleted{\unsure{donc}} L'anneau affine de $\underline{\mathrm{End}}_k(M)$ et a fortiori celui de $\mathrm{GL}(n)$ s'en va \uline{sur} $A$, OK.
% archive sheet number: 30
